% =====================================================================
%  PARTIE 5 — Limites, défis et perspectives
% =====================================================================
\chapter{Limites, défis et perspectives d'avenir}
\label{ch:limites}

\section{Des ressources matérielles limitées}

La principale difficulté du TinyML consiste à faire fonctionner un modèle utile
sur un appareil disposant de ressources réduites. La mémoire vive, l'espace de
stockage, la puissance de calcul et l'énergie disponible limitent la complexité
des modèles déployables. L'optimisation doit ainsi prendre en compte plusieurs
critères simultanément : la taille du modèle, sa précision, son temps
d'exécution et sa consommation~\cite{banbury2020benchmarking}.

Les outils logiciels imposent également des contraintes. TensorFlow Lite for
Microcontrollers, par exemple, prend en charge un ensemble limité d'opérations
et nécessite une gestion adaptée de la mémoire~\cite{david2021tflm}. Un modèle
conçu sur ordinateur ne peut donc pas toujours être transféré directement sur un
microcontrôleur sans modification.

Ces contraintes conduisent à rechercher un compromis : le modèle le plus
complexe n'est pas nécessairement le plus pertinent si un modèle plus simple
répond au besoin tout en consommant moins de ressources.

\section{Une fiabilité dépendante des données et du contexte}

Les performances d'un modèle dépendent des données utilisées pour son
entraînement et de leur représentativité. Dans les applications vocales, le
bruit peut dégrader la reconnaissance. Dans les dispositifs portables, les
mouvements, les différences entre utilisateurs ou les variations de
positionnement des capteurs peuvent modifier les signaux analysés.

Un modèle peut donc donner de bons résultats pendant les essais puis devenir
moins fiable lorsque les conditions changent. Il est nécessaire de le tester
dans des situations variées et de suivre ses performances après le
déploiement. Cette exigence est particulièrement importante lorsque les erreurs
peuvent avoir des conséquences sur la surveillance d'un équipement ou sur
l'interprétation de données de santé.

\section{L'entraînement et l'évolution des modèles}

Dans de nombreux systèmes TinyML, l'entraînement est effectué sur un ordinateur
ou une infrastructure distante, tandis que le microcontrôleur réalise uniquement
l'inférence, c'est-à-dire l'utilisation du modèle déjà entraîné. Cette
séparation s'explique par les besoins généralement plus élevés de
l'entraînement en mémoire et en puissance de calcul. TensorFlow Lite for
Microcontrollers ne prend d'ailleurs pas en charge l'entraînement sur
l'appareil~\cite{tflm-doc}.

Cette organisation soulève la question de l'évolution du système. Lorsque les
données ou les conditions d'utilisation changent, une nouvelle version du modèle
peut devenir nécessaire ; elle doit alors être redéployée, par exemple par une
mise à jour du firmware à distance (OTA).

L'apprentissage directement sur l'appareil constitue une piste d'évolution, des
travaux ayant montré qu'il est possible d'entraîner partiellement un réseau
avec moins de 256~Ko de mémoire~\cite{lin2022ondevice}. Il n'est cependant pas
encore une solution générale pour tous les microcontrôleurs : il reste limité
par les ressources disponibles et par la complexité supplémentaire qu'il
introduit.

\section{La sécurité des modèles embarqués}

Placer un modèle dans un objet physique l'expose à de nouvelles menaces : un
attaquant ayant accès à l'appareil peut tenter d'extraire le modèle de la
mémoire flash, de le modifier ou de fournir au capteur des entrées conçues pour
le tromper. La protection de la mémoire (chiffrement de la flash, démarrage
sécurisé) et la sécurisation des mises à jour OTA, que proposent par exemple
les ESP32~\cite{espressif-esp32}, font donc partie intégrante de la conception
d'un système AIoT.

\section{Vers des architectures complémentaires entre l'objet et le cloud}

Le développement du TinyML ne signifie pas nécessairement la disparition du
cloud. Les deux approches peuvent remplir des rôles complémentaires : l'objet
exécute les analyses locales compatibles avec ses ressources, tandis qu'une
infrastructure plus puissante peut prendre en charge l'entraînement et les
traitements plus exigeants.

\begin{figure}[H]
\centering
\begin{tikzpicture}[node distance=2.4cm]
  \node[blocv, text width=2.9cm] (o) {Objets TinyML\\\footnotesize inférence locale, alertes};
  \node[bloc, text width=2.9cm, right=of o] (e) {Passerelle / edge\\\footnotesize agrégation, analyses};
  \node[bloco, text width=2.9cm, right=of e] (c) {Cloud\\\footnotesize entraînement, supervision};
  \draw[fleche] ([yshift=4pt]o.east) -- node[above,font=\scriptsize]{résultats} ([yshift=4pt]e.west);
  \draw[fleche] ([yshift=4pt]e.east) -- node[above,font=\scriptsize]{données utiles} ([yshift=4pt]c.west);
  \draw[fleche, dashed] ([yshift=-4pt]c.west) -- node[below,font=\scriptsize]{nouveau modèle} ([yshift=-4pt]e.east);
  \draw[fleche, dashed] ([yshift=-4pt]e.west) -- node[below,font=\scriptsize]{mise à jour OTA} ([yshift=-4pt]o.east);
\end{tikzpicture}
\caption{Répartition des traitements dans une architecture AIoT hybride.}
\label{fig:hybride}
\end{figure}

Les perspectives reposent notamment sur l'amélioration de l'efficacité des
modèles et de leur adaptation aux contraintes matérielles. Les travaux récents
sur la quantification montrent que la recherche vise à réduire les besoins en
mémoire et en énergie tout en préservant des performances adaptées à
l'application~\cite{gholami2021survey}. Parallèlement, les microcontrôleurs
intègrent de plus en plus d'accélérateurs neuronaux dédiés, comme le montre le
STM32N6 présenté dans la partie~\ref{ch:materiel}.

L'enjeu n'est donc pas de placer toute l'intelligence dans chaque capteur, mais
de répartir les traitements de façon pertinente selon les besoins de
réactivité, d'autonomie et de confidentialité.
